\section{Weekly Summary}

\begin{tcolorbox}[colback=yellow!10,colframe=black!50,title=AI-Generated Content]
\noindent The summary below was written by Claude (Anthropic) from that week's regression outputs and series descriptions. It has not been reviewed or fact-checked by a human, and should be read as an automated, informal recap -- not analysis.
\end{tcolorbox}

This week the cosmic dice of FRED rolled us a truly inspired matchup: Labor Compensation for Utilities: Electric Power Generation, Transmission and Distribution in the United States, pitted valiantly against the Economic Policy Uncertainty Index for South Korea. Yes, that South Korea. And yes, that index has been discontinued, which means we are essentially performing statistical seances on a dead data series. Nothing says rigorous science like regressing American electrician paychecks against the vibes of Korean policymakers who stopped reporting their anxiety in 2014.

And yet -- brace yourselves -- the numbers came back looking suspiciously respectable. We got an R-squared of about 0.41, a coefficient of roughly 2.28 with a p-value of 0.0006, and an F-test that also clocked in at 0.0006. In the sacred temple of "statistical significance," these are the kinds of numbers people frame and hang on their walls. My hands are trembling. Could it be that when South Koreans felt uncertain about economic policy, American power grid workers somehow got paid more? Truly the invisible hand works in mysterious, transcontinental, chronologically-mismatched ways.

Now, before anyone rushes off to build a hedge fund around "Korean Angst Utility Wage Arbitrage," let me gently remind you that this entire project is two randomly grabbed time series shoved into a regression like reluctant guests at a dinner party. One is annual, one is monthly. One ends in 2025, the other gave up on life in 2014. There is no causal story here, no mechanism, no theory -- just two lines that happen to drift upward together over time, which is roughly as meaningful as noticing that both my houseplant and the national debt have grown since 2020. Time trends are the great matchmaker of spurious correlation, and this week they earned their commission.

So enjoy the deeply impressive-looking p-value for exactly what it is: a delightful coincidence produced by a script that has no idea what electricity or Korea are. As always, the moral of the story is that correlation is cheap, randomness is generous, and if you squint hard enough, everything correlates with everything. See you next week for another beautiful, meaningless graph.
